\documentclass{standalone}
\usepackage{ctex}
\usepackage{tikz}

\usetikzlibrary{cartes}

% Styles
\tikzstyle{axes}=[]
\tikzstyle{important line}=[very thick]
\tikzstyle{information text}=[rounded corners,fill=red!10,inner sep=1ex]

\begin{document}

\begin{tikzpicture}
  \tikzmath{
    \A{1} = 0; \A{2} = 3; 
    \B{1} = -1; \B{2} = -2;
    \C{1} = 3; \C{2} = -2;
    \kd = -0.5;
    \ke = 0.3;
  }

  \tikzset{
    plane/combine={\B,\C,\kd,1-\kd,\D},
    plane/combine={\A,\C,\ke,1-\ke,\E},
    plane/intersect ll={\A,\B,\D,\E,\F},
    plane/circumcenter={\A,\B,\C,\Oa},
    plane/length={\Oa,\A,\ra},
    plane/circumcenter={\A,\E,\F,\Ob},
    plane/length={\Ob,\A,\rb},
    plane/circumcenter={\D,\B,\F,\Oc},
    plane/length={\Oc,\D,\rc},
    plane/circumcenter={\D,\C,\E,\Od},
    plane/length={\Od,\D,\rd},
    plane/intersect cc={\Oa,\ra,\Ob,\rb,\P,\Q},
    plane/exclude={\P,\Q,\A,\M},
    plane/project={\A,\M,\B,\Pa},
    plane/project={\B,\M,\C,\Pb},
    plane/project={\C,\M,\A,\Pc},
    plane/project={\D,\M,\E,\Pd},
    plane/orthocenter={\A,\B,\C,\Ha},
    plane/orthocenter={\A,\E,\F,\Hb},
    plane/orthocenter={\D,\B,\F,\Hc},
    plane/orthocenter={\D,\C,\E,\Hd},
  }

  % \axes[grid] (-4:8,-4:4);
  \coordinate (A) at (\A{1},\A{2});
  \coordinate (B) at (\B{1},\B{2});
  \coordinate (C) at (\C{1},\C{2});
  \coordinate (D) at (\D{1},\D{2});
  \coordinate (E) at (\E{1},\E{2});
  \coordinate (F) at (\F{1},\F{2});
  \coordinate (M) at (\M{1},\M{2});
  \coordinate (O1) at (\Oa{1},\Oa{2});
  \coordinate (O2) at (\Ob{1},\Ob{2});
  \coordinate (O3) at (\Oc{1},\Oc{2});
  \coordinate (O4) at (\Od{1},\Od{2});
  \coordinate (P1) at (\Pa{1},\Pa{2});
  \coordinate (P2) at (\Pb{1},\Pb{2});
  \coordinate (P3) at (\Pc{1},\Pc{2});
  \coordinate (P4) at (\Pd{1},\Pd{2});
  \coordinate (H1) at (\Ha{1},\Ha{2});
  \coordinate (H2) at (\Hb{1},\Hb{2});
  \coordinate (H3) at (\Hc{1},\Hc{2});
  \coordinate (H4) at (\Hd{1},\Hd{2});

  \draw[thick,navy] (A) -- (B) -- (D) -- (F) (A) -- (C);
  \draw[thick,red] ($(P1)!-0.25!(P2)$) -- ($(P1)!1.25!(P2)$);
  \draw[thick,violet] ($(H2)!-0.25!(H4)$) -- ($(H2)!1.25!(H4)$);
  % \draw[teal] (O1) circle (\ra);
  % \draw[violet] (O2) circle (\rb);
  % \draw[purple] (O3) circle (\rc);
  % \draw[magenta] (O4) circle (\rd);
  
  \foreach \p/\placement in {A/above,B/below left,C/below right,D/below right,E/below left,F/above left,M/above right}{
    \fill[red] (\p) circle (1.5pt);
    \draw (\p) node[\placement] {$\p$};
  }

  \foreach \p/\placement in {P1/above,P2/above,P3/above,P4/above}{
    \fill[orange] (\p) circle (1.5pt);
  }

  \foreach \p/\placement in {H1/above,H2/above,H3/above,H4/above}{
    \fill[orange] (\p) circle (1.5pt);
  }

  \draw[xshift=-2.5cm,yshift=-5.5cm] node [right,text width=9cm,style=information text]
  {
    \textit{四个三角形的垂心共线，记为 Gauss 线。Gauss 线与 Steiner 线平行，且点 M 到 Gauss 线的距离是到 Steiner 线的 2 倍。}
  };

\end{tikzpicture}

\end{document}